%! The Fundamental Theorem of Linear Algebra — Unit 3, Lesson 3.6 — Warm-Up
\documentclass[10pt]{article}
\usepackage{linalg-article}
\usepackage{linalg-boxes}
\newcommand{\TallMath}[1]{$\displaystyle #1\rule[-1.4em]{0pt}{3.2em}$}
\newcommand{\xx}{\mathbf{x}}
\newcommand{\vv}{\mathbf{v}}
\newcommand{\ww}{\mathbf{w}}
\newcommand{\zero}{\mathbf{0}}
\newcommand{\T}{\mathsf{T}}

\begin{document}

\pageheader{Unit 3, Lesson 3.6}{Warm-Up}
\namedateperiod

\vspace{0.12in}

\begin{notesbox}{Getting Ready: Skills We'll Use Today}
\small These skills from Lessons 1.2, 1.3, and 3.5 power today's capstone. Answer quickly.

\vspace{4pt}
\textbf{1. Perpendicular $\Leftrightarrow$ dot $=0$ (Lesson 1.2).} Two vectors are perpendicular exactly when
their dot product is $0$. Compute each dot product and circle perpendicular or not.
\[
  \vv=\begin{bsmallmatrix} 1\\0\\1 \end{bsmallmatrix},\ \ww=\begin{bsmallmatrix} -1\\-1\\1 \end{bsmallmatrix}:
  \quad \vv\cdot\ww=\blank{1.4cm}\ \ \text{(perp / not)}
  \qquad
  \vv=\begin{bsmallmatrix} 1\\2\\3 \end{bsmallmatrix},\ \ww=\begin{bsmallmatrix} 1\\1\\1 \end{bsmallmatrix}:
  \quad \vv\cdot\ww=\blank{1.4cm}\ \ \text{(perp / not)}
\]

\vspace{4pt}
\textbf{2. $A\xx$ row by row (Lesson 1.3).} Each entry of $A\xx$ is a \emph{row of $A$ dotted with $\xx$}.
Fill in the two dot products, then $A\xx$.
\[
  A=\begin{bmatrix} 1 & 0 & 1 \\ 0 & 1 & 1 \end{bmatrix},\quad
  \xx=\begin{bmatrix} -1 \\ -1 \\ 1 \end{bmatrix}
  \qquad
  A\xx=\begin{bmatrix} \text{row}_1\cdot\xx \\[3pt] \text{row}_2\cdot\xx \end{bmatrix}
      =\begin{bmatrix}\rule{0pt}{1.1em}\blank{0.8cm}\\[3pt] \blank{0.8cm}\end{bmatrix}
\]

\vspace{4pt}
\textbf{3. Four dimensions from $r$ (Lesson 3.5).} A matrix $A$ is $3\times 5$ with rank $r=2$. Give the four
dimensions.
\[
  \dim C(A)=\blank{0.9cm}\quad \dim C(A^{\T})=\blank{0.9cm}\quad
  \dim N(A)=\blank{0.9cm}\quad \dim N(A^{\T})=\blank{0.9cm}
\]
\end{notesbox}

\end{document}
